% TOP_PROBLEM: {"id":30007599,"problem_number":"LOCAL-30007599","title":"Core-market fragility","category_id":21,"category":{"id":21,"name":"economics","display_name":"Economics","description":"Open economic theory statements, published research agendas, and proposed scoped specifications; question provenance and historical priority are qualified per record.","slug":"economics","order_index":21,"created_at":"2026-10-08T00:00:00Z"},"status":"research_question","external_url":null,"legacy_ids":[],"published":false,"statement":"In the explicitly specified setting: What trading, collateral and settlement arrangements keep government-bond and other core markets liquid during extreme shocks? The mathematical statement above fixes the target, assumptions and scope of this catalogue formulation.","economics":{"original_id":88,"field":"Asset pricing and financial markets","kind":"design","formalization_basis":"adapted_research_question","source_status":"research_agenda","priority_status":"earliest_found","priority_note":"This is the earliest explicit statement verified in this audit; first priority for the full catalogue formulation is not established.","earliest_verified_reference":{"authors":["Bank of England"],"title":"Bank of England Agenda for Research, 2025–2028","year":2026,"url":"https://www.bankofengland.co.uk/research/bank-of-england-agenda-for-research","doi":null,"locator":"The Financial System / Financial market dynamics, paragraphs 3–4; Priority topics for 2026, item 1","evidence_summary":"Questions clearing, netting and core-market resilience."},"current_reference":{"authors":["Sirio Aramonte","Andreas Schrimpf","Hyun Song Shin"],"title":"Non-bank financial intermediaries and financial stability","year":2021,"url":"https://www.bis.org/publications/working-paper-972-non-bank-financial-intermediaries-and-financial-stability.pdf","doi":null,"locator":"Section 5, Open questions for policy and research, printed pp.38–39 (PDF pp.41–42), January 2022 revision","evidence_summary":"Explicitly asks optimal oversight and public-backstop design."},"formalization_note":"Adapts an explicitly verified research-agenda passage. Mathematical objects, interventions, objectives and scope are proposed here; existing model-specific solutions are not relabeled unsolved."},"statusQualification":"Published question or research agenda verified at the cited locator. The mathematical specification is adapted for this catalogue and includes additional assumptions; source publication does not establish first-ever priority or certify this exact model remains unsolved. Adapts an explicitly verified research-agenda passage. Mathematical objects, interventions, objectives and scope are proposed here; existing model-specific solutions are not relabeled unsolved.","statusReviewedAt":"2026-10-08"}
% FIRST_POSTED: 2026-10-08
% ENTRY_KIND: statement_only
% !TeX program = lualatex
\documentclass[11pt]{article}
\usepackage{fontspec}
\usepackage[a4paper,margin=25mm]{geometry}
\usepackage{amsmath,amssymb,mathtools}
\usepackage{xurl}
\usepackage[hidelinks]{hyperref}
\newcommand{\sref}[2]{\hyperlink{source:#1:#2}{[S#2]}}
\setlength{\emergencystretch}{3em}
\begin{document}
% problemId:30007599
\section*{Core-market fragility}\label{30007599}
\subsection{Definitions and mathematical statement}
Fix a government-bond market with dealers, leveraged funds, customer orders, collateral and a specified settlement network. A feasible architecture \(a\in A\) specifies trading access, central clearing, netting periods, collateral eligibility and margin schedules. At each date agents choose portfolios and liquidity buffers optimally, taking these rules into account. Let \(L(a,\omega)\) be realized social loss in shock state \(\omega\), combining execution costs, failed settlement, default externalities and real-economy disruption, with its components separately measured. Let \(C(a)\) be the ordinary resource and compliance cost. For a publicly defined uncertainty set \(\mathcal P\) of shock distributions, seek
\[a^*\in\arg\min_{a\in A}\{C(a)+\sup_{P\in\mathcal P}E_P[L(a,\omega)]\}.\]
The research task is to construct an empirically disciplined loss function and characterize designs that remain resilient when normally liquid assets experience simultaneous selling and funding shocks. Specify institutional budget, legal and incentive constraints defining \(A\). Account for changes in leverage and liquidity provision induced by the architecture, rather than holding agent behavior fixed. Report the frontier between routine intermediation costs and tail losses, and the robustness of optimal designs to uncertain market-impact elasticities and shock dependence. A design optimal under one calibrated stress episode need not solve the wider market-resilience question.

\par\textbf{Formulation and scope.} Adapts an explicitly verified research-agenda passage. Mathematical objects, interventions, objectives and scope are proposed here; existing model-specific solutions are not relabeled unsolved.

\par\textbf{Open-question provenance.} An explicit question or research agenda was verified at the source locator. This does not by itself certify that every later version remains unresolved \sref{30007599}{1}.

\par\textbf{Historical priority.} This is the earliest explicit statement verified in this audit; first priority for the full catalogue formulation is not established.

\subsection{Short English statement}
In the explicitly specified setting: What trading, collateral and settlement arrangements keep government-bond and other core markets liquid during extreme shocks? The mathematical statement above fixes the target, assumptions and scope of this catalogue formulation.

\subsection{Sources}
\begin{itemize}\raggedright
\item[\textnormal{[S1]}]\hypertarget{source:30007599:1}{} Bank of England. 2026. Bank of England Agenda for Research, 2025–2028. The Financial System / Financial market dynamics, paragraphs 3–4; Priority topics for 2026, item 1.
\url{https://www.bankofengland.co.uk/research/bank-of-england-agenda-for-research}.
{\small\textit{Earliest verified question/agenda reference; historical priority qualified below. Questions clearing, netting and core-market resilience.}}

\item[\textnormal{[S2]}]\hypertarget{source:30007599:2}{} Sirio Aramonte, Andreas Schrimpf, Hyun Song Shin. 2021. Non-bank financial intermediaries and financial stability. Section 5, Open questions for policy and research, printed pp.38–39 (PDF pp.41–42), January 2022 revision.
\url{https://www.bis.org/publications/working-paper-972-non-bank-financial-intermediaries-and-financial-stability.pdf}.
{\small\textit{Supporting or current literature; see evidence qualification. Explicitly asks optimal oversight and public-backstop design.}}
\end{itemize}
\end{document}
